\section{Lean formalization plan}
\label{sec:lean-plan}

This section gives a concrete path from the mathematical framework to a complete Lean
formalization in the repository.  The target environment is Lean 4.32.0 \cite{lean4} with the
matching mathlib revision, as pinned by the current ProveIt root.  The intended library
location is
\[
\code{Logic/PrimitiveRecursiveArithmetic/Derivations/Lean/IPRADerivations/FRNF/}.
\]

\subsection{Reuse before abstraction}

The implementation should reuse the checked modules rather than create parallel
arithmetic:

\begin{center}
\begin{tabularx}{\textwidth}{@{}p{0.34\textwidth}X@{}}
\toprule
Existing module & Reuse \\
\midrule
\code{Core/Basic} & \code{inst}, \code{ind}, defining equations, congruence. \\
\code{Core/Semiring}, \code{Core/Order}, \code{Core/Monus} & Boolean arithmetic,
max/min, inequalities, error budgets. \\
\code{Core/Pairing} & Pair codes, projections, equality test, simultaneous recursion. \\
\code{Core/Search} & Least bounded search and minimality. \\
\code{Core/BoundedSums}, \code{Core/MultiSums} & Booleanized bounded quantifiers and
finite folds. \\
\code{Core/Definitions}, \code{Core/Sums} & Generated primitive-recursive recurrences. \\
\code{Core/Descent}, \code{Core/DescentEx} & Generic finite/infinite descent endpoints. \\
\code{Core/Division}, \code{Core/Gcd}, \code{Core/Bezout} & Majorized Euclidean examples
and denominator arithmetic. \\
\code{NativeBiInterpretation} & Final \code{pra_iff_native} export. \\
\bottomrule
\end{tabularx}
\end{center}

\subsection{Proposed module sequence}

\paragraph{1. \code{FRNF/BoundedFormula.lean}.}
Define the bounded formula syntax independently of host propositions:

\begin{lstlisting}
inductive BForm : Nat -> Type
  | eq    (s t : PTm n)
  | bot
  | and   (A B : BForm n)
  | or    (A B : BForm n)
  | imp   (A B : BForm n)
  | allLt (bound : PTm n) (body : BForm (n+1))
  | exLt  (bound : PTm n) (body : BForm (n+1))
\end{lstlisting}

Implement \code{toFm}, generated characteristic symbols, and substitution.  Prove
simultaneously:

\begin{lstlisting}
theorem char_boolean (A : BForm n) :
  [] |-[(PRAAx)] (charTm A =' 0) \/ (charTm A =' 1)

theorem to_char (A : BForm n) :
  [A.toFm] |-[(PRAAx)] charTm A =' 1

theorem of_char (A : BForm n) :
  [charTm A =' 1] |-[(PRAAx)] A.toFm
\end{lstlisting}

The bounded quantifier cases should invoke generic bounded sums/products.  Prove
\code{char_subst} early; otherwise every later generated invariant becomes painful.

\paragraph{2. \code{FRNF/Table.lean}.}
Choose one list/array coding based on pairing and simultaneous recursion.  Define
\code{len}, \code{get}, \code{set}, \code{append}, \code{foldMax}, and a bounded map.
Prove object equations for in-range lookup, update, and fold invariants.  Avoid a
universal function argument: instantiate folds with fixed bodies or a fixed finite
opcode family.

\paragraph{3. \code{FRNF/Horizon.lean}.}
Represent each concrete horizon by a \code{PR} symbol.  Define prefix maximum and prove
monotonicity, domination, and dual-horizon composition.  A useful object record is a
metatheoretic structure whose fields include the symbols and the embedded derivations:

\begin{lstlisting}
structure DualHorizonCert (r j : PR 1) where
  read_le_justify : [] |-[(PRAAx)]
    allOpen (app r ![var 0] <=' app j ![var 0])
  read_mono : ...
  justify_mono : ...
\end{lstlisting}

The structure itself is host data, but every field used by the compiler is an embedded
proof.

\paragraph{4. \code{FRNF/Demand.lean}.}
For a fixed generated graph, first formalize a specialized demand pass.  Then factor a
generic encoded-graph theorem if worthwhile.  The main endpoint is a QF loop proving
that the reverse array after \(m\) steps is closed and least.  Separate read and proof
arrays in the state.

\paragraph{5. \code{FRNF/Trace.lean}.}
Move the source-level \code{QFLoop} and proof-schedule combinators from this archive
into the checked library.  Add a record for a concrete transition with:
\code{state}, \code{init}, \code{step}, \code{fuel}, \code{inv}, and embedded base/step
proofs.  The compiler theorem should return \code{Prf PRAAx}, not a host predicate.

\paragraph{6. \code{FRNF/Viability.lean}.}
Encode level offsets, candidate counts, validity bits, edge bits, and viability masks.
Implement the backward dynamic program and forward least path.  Prove:

\begin{lstlisting}
theorem viable_iff_suffix_path : ...
theorem selected_path_valid : ...
theorem selected_path_compatible : ...
\end{lstlisting}

The first theorem is most easily proved by QF induction on the reversed level index;
existence of a suffix path is represented by a bounded Boolean characteristic, not an
unbounded existential in the invariant.

\paragraph{7. \code{FRNF/Stability.lean}.}
Define finite diameter, stable observers, gap collapse, and the shrinking-viability
one-step estimate.  Initially specialize to rational observables encoded by signed
integer pairs.  General metric presentations can be layered later.

\paragraph{8. \code{FRNF/FiniteNet.lean}.}
Define a finite rational net interface and exact minimum fold.  Prove the finite
comparison contract from explicit coverage and modulus formulas.  Keep the semantic
real infimum theorem in a separate host module; the object endpoint should quantify
only over coded inputs and witnesses.

\paragraph{9. \code{FRNF/Guarded.lean}.}
Formalize rational absolute values, positive denominators, and the cleared-denominator
version of \cref{thm:guarded-reciprocal}.  Export reusable precision transformers for
reciprocal, quotient, and normalized vectors.

\paragraph{10. \code{FRNF/Compiler.lean}.}
Define the proof schedule over formulas and a generated node descriptor.  A fully
generic reflection layer is optional.  It is sufficient that a metaprogram produces
concrete Lean definitions and applies \code{Schedule.extend}.  Add final native export.

\subsection{Formula generation versus opaque declarations}

The roadmap file in this package uses \code{opaque} formulas to make endpoints
unambiguous.  The completed implementation must replace every opaque endpoint with an
explicit \(\Fm\) expression assembled from the actual PR symbols.  Otherwise a proof
of the opaque formula would not establish the intended theorem.

A useful discipline is to define each object formula immediately next to a human
pretty-printer and a host evaluation theorem.  The pretty-printer assists review; the
host theorem assists debugging; only the explicit object formula and its embedded
proof establish the \(\iPRA\) result.

\subsection{Suggested first end-to-end milestone}

The smallest convincing milestone is not the full generic graph compiler.  It is a
single integrated theorem:

\begin{quote}
Given a fixed two-stage transducer with distinct computation and justification
horizons, compile the reverse demand bounds, materialize a finite table, verify one
clocked loop, and return a closed \(\Prf\,\PRAAx\) derivation; then export it through
\code{pra_iff_native}.
\end{quote}

This exercises every architectural boundary while keeping domain arithmetic trivial.
The second milestone should add a finite relational cloud and stable observer.

\subsection{Build and audit procedure}

A result should be labeled object-checked only after the following commands succeed in
the pinned repository environment:

\begin{lstlisting}[language=bash]
lake build IPRADerivations
lake env lean Logic/PrimitiveRecursiveArithmetic/Derivations/Lean/\
  IPRADerivations/FRNF/Audit.lean
\end{lstlisting}

The audit module should contain \code{#print axioms} for every exported theorem and a
machine-readable list of declarations containing \code{sorryAx}.  The acceptable
metatheoretic axioms should match the established repository policy; no result with
\code{sorryAx} is object-checked.

\subsection{What the accompanying Lean files establish}

\code{lean/FRNF/Kernel.lean} contains the logical schedule, QF-loop closure, stable
observer, boundary discharge, gap collapse, and native export using the inspected
repository API.  It contains no \code{axiom} declaration and no \code{sorry}, but it
was not compiled in this environment.  \code{lean/FRNF/ObjectTargets.lean} deliberately
contains visible \code{sorry} markers for Booleanization, demand closure, viability,
stable clouds, finite nets, and guarded reciprocals.  These are a roadmap, not checked
results.  This status split should be preserved during migration.
